% ======================================================================
%  41_layer3_notes.tex — 外部注释卡（L3 注释层）
%  依赖 preamble.tex；独立自检目录 build/check_notes/
% ======================================================================

\section{外部注释卡：读原文真正会卡住的概念}

这一章不是术语表，也不是百科。它只做一件事：当原文第一次把某个概念当作工具使用时，给读者一个可以独立读完的入口，读完就能回到原文那一节继续往下读。每张卡都是同一个五段格式：

\begin{enumerate}[label=\arabic*.,leftmargin=1.8em,itemsep=1pt]
  \item \textbf{一句话}：用一句人话说清它是什么，以及为什么非有它不可；
  \item \textbf{数学要点}：三到六句，定义、关键性质、为什么重要；
  \item \textbf{在 K--W 中的位置}：它出现在哪一节、对应骨架节点 S?/N?、被用来做什么；
  \item \textbf{读原文前需要的前置}：一到三张更基础的卡的编号；
  \item \textbf{常见坑}：一两条最容易把人带偏的误解。
\end{enumerate}

每张卡末尾给出所据条目的完整链接。所有链接均经实际抓取核对（HTTP 200，且页面内容与本卡陈述相符），抓取日期 2026-11-15；若某个 nLab 条目日后改名，按卡内给出的条目名在 nLab 站内检索即可。链接以 nLab 为主，因为本机访问 nLab 稳定；Wikipedia 的数学条目在中文网络环境下不可靠，故不采用。\anchor{src: nLab 在线条目，2026-11-15 抓取}

\textbf{怎么排的。}卡片按“读者读论文的顺序”而非按字母排：第 1--13 张是数学前置（丛、模空间、超凯勒、Higgs、Hitchin、非阿贝尔 Hodge、广义复几何、平坦丛、Riemann--Hilbert），第 14--24 张是物理语言（A/B 模型、镜像、膜、扭化、$N=4$、BPS、BRST、瞬子、单极子、Bogomolny 以及 $SL(2,\ZZ)$ 与模形式），第 25--35 张是后半程的代数几何机器（相干层、导出范畴、Fourier--Mukai、D-模、仿射 Grassmann、Hecke、几何 Satake、几何 Langlands、对偶群、根格），第 36 张供附录对照（顶点算子代数与 CFT 路线）。\anchor{src: plan/04 骨架编号表}

\textbf{怎么读。}完整路线建议按 1--13 顺读；若只有半天，最小闭环是第 3、4、6、7、8 张（超凯勒、超凯勒商、Hitchin 模空间、Hitchin 方程、Hitchin 纤维化）加上第 9 张（非阿贝尔 Hodge），这六张卡撑起 §4 之后的全部几何舞台。读 §5 之前补第 10 张（广义复几何）；读 §8--§9 之前补第 11、15、16、30、31 张；读 §11 之前补第 26、28 张。每张卡“前置”一栏给出严格依赖，不必按编号顺序全部预习。\anchor{src: plan/01 分层讲义规范}

\subsection*{第一组\quad 数学前置：丛、模空间、Hitchin 几何}

\begin{kcard}{注释 1：主丛与稳定丛（principal bundle and stable bundle）}
\anchor{nLab: stable+bundle}

\textbf{一句话}：规范场住在主丛上，而“稳定”是把无穷多个同构类压成一个好模空间的筛选条件。

\textbf{数学要点}：$G$-主丛 $P\to X$ 是纤维为 $G$、$G$ 在每根纤维上自由且传递作用的丛；联络 $A\in\Omega^1(P,\mathfrak{g})$ 给出曲率 $F=\dd A+A\wedge A$，规范场就是联络在丛上的局部实现。取 $G=GL(n)$（或 $U(n)$）时，主丛的关联丛就是秩 $n$ 的向量丛，全纯主丛则给出全纯向量丛及其模空间 $\mathrm{Bun}_G(X)$。在紧 Riemann 面 $X$ 上，向量丛 $E$ 的斜率是 $\mu(E)=\deg(E)/\mathrm{rk}(E)$，\term{stable}{稳定}指每个真子丛 $F\subset E$ 都满足 $\mu(F)<\mu(E)$（半稳定把 $<$ 换成 $\le$）。Mumford 的几何不变量理论把稳定点整理成商，而 Narasimhan--Seshadri 定理说 $\deg=0$ 的稳定全纯丛恰好对应不可约平坦联络。

\textbf{在 K--W 中的位置}：§4 取 $X=\Sigma$ 为紧 Riemann 面做紧化，$\mathrm{Bun}_G(\Sigma)$ 与以它为底空间的 $M_H$ 依次登场（N4、S3）；§9 的 Hecke 算子作用在 $\mathrm{Bun}_G$ 上（N9、S7）。后文所有“电特征膜”“Hecke 特征层”都以 $\mathrm{Bun}_G$ 或 $M_H$ 上的对象身份出现。

\textbf{读原文前需要的前置}：无，本卡是整组的地基。

\textbf{常见坑}：稳定不等于半稳定；半稳定边界让模空间只能以 GIT 商而不是轨道空间的形式存在。稳定性还依赖极化的选择（$\deg$ 由 Kähler 形式决定），换 Kähler 类可能换掉稳定集合。

\noindent\textbf{链接}：\url{https://ncatlab.org/nlab/show/stable+bundle}
\end{kcard}

\begin{kcard}{注释 2：模空间与 GIT（moduli space and geometric invariant theory）}
\anchor{nLab: geometric+invariant+theory}

\textbf{一句话}：GIT 是“先删掉不稳定点、再取商”的标准机器，模空间由此重新成为一个代数簇。

\textbf{数学要点}：设约化群 $G$ 作用在射影簇 $X$ 上，$X^{ss}$ 为半稳定点集，则 GIT 商写成 $X^{ss}/\!/G=\mathrm{Proj}\bigoplus_k H^0(X,L^{\otimes k})^G$；在稳定点集 $X^{s}$ 上这个商是几何商，局部是轨道空间（orbifold 或 Deligne--Mumford 栈）。经典应用是“模空间 of bundles”：规范群作用在参数空间上，稳定点恰是斜率稳定丛，商即模空间。对 Higgs 丛，把 $G$ 换成规范群、把参数空间换成 Higgs 对的空间，同一机器给出 $M_H$ 的代数构造。GIT 的意义在于：模空间上的函数就是规范不变量，于是“在模空间上做几何”才有意义。

\textbf{在 K--W 中的位置}：§4 用模空间语言定义 $M_H$（N4、S3）；§11 的 B-模型把相干层放在 $M_H$ 上，其“良好性”正依赖 GIT 商的存在（N11）。

\textbf{读原文前需要的前置}：注释 1。

\textbf{常见坑}：$X^{ss}/\!/G$ 不是普通轨道空间，轨道可能被合并；GIT 稳定与 Hilbert--Mumford 稳定、斜率稳定分属不同层级，不能混用；非约化群上这套机器不适用。

\noindent\textbf{链接}：\url{https://ncatlab.org/nlab/show/geometric+invariant+theory}
\end{kcard}

\begin{kcard}{注释 3：超凯勒流形（hyperkähler manifold）}
\anchor{nLab: hyperkähler+manifold}

\textbf{一句话}：超凯勒流形是一个流形上同时装备三个两两反对易的复结构，而且三个复结构都 Kähler。

\textbf{数学要点}：$4n$ 维黎曼流形称为超凯勒的，若其和乐群含于 $\mathrm{Sp}(n)$；等价地，它带有三个复结构 $I,J,K$，满足 $I^2=J^2=K^2=IJK=-1$，并且相应的三个 Kähler 形式 $\omega_I,\omega_J,\omega_K$ 都闭。任意单位向量 $aI+bJ+cK$ 仍是复结构，因此复结构构成一个 $\mathbb{CP}^1$ 族，这是超凯勒几何最实用的性质。紧致例子包括 K3 面、K3 面或 4-环面的 Hilbert scheme（Beauville 构造），以及 O'Grady 的两个例外例子；物理侧，$D=3$、$N=4$ 超杨--米尔斯理论的 Coulomb 分支与 Higgs 分支是超凯勒的（Seiberg--Witten）。

\textbf{在 K--W 中的位置}：§4 指出 $M_H$ 是超凯勒流形（N4、S3），这是全文的几何基石；三个复结构决定了膜的“类型”：B-模型需要复结构，A-模型需要辛结构，而 $M_H$ 上不同复结构的选取对应不同 twist（§5，N5；§11--§12）。后文所谓“$I$、$J$、$K$ 中的哪一个”不是装饰，而是决定该膜是 A-型还是 B-型的关键。

\textbf{读原文前需要的前置}：注释 1。

\textbf{常见坑}：超凯勒（和乐 $\mathrm{Sp}(n)$）不等于四元数 Kähler（和乐 $\mathrm{Sp}(n)\mathrm{Sp}(1)$）；“三个复结构”也不是任取三个复结构，必须满足四元数关系且同时 Kähler。

\noindent\textbf{链接}：\url{https://ncatlab.org/nlab/show/hyperk%C3%A4hler+manifold}
\end{kcard}

\begin{kcard}{注释 4：超凯勒商（hyperkähler quotient）}
\anchor{nLab: symplectic+reduction}

\textbf{一句话}：超凯勒商是把辛商（Marsden--Weinstein 商）同时对三个辛形式各做一遍，得到的商空间仍然超凯勒。

\textbf{数学要点}：设超凯勒流形 $M$ 上有保持超凯勒结构的群 $G$ 作用，超凯勒矩映射是三分量映射 $\mu=(\mu_I,\mu_J,\mu_K):M\to\mathfrak{g}^*\otimes\mathbb{R}^3$；超凯勒商定义为 $M/\!/\!/G:=\mu^{-1}(0)/G$，它继承 $M$ 的超凯勒结构。在任意一个复结构下看，它就是普通的辛商；三个方程 $\mu_I=\mu_J=\mu_K=0$ 合起来是同一组实方程的三个分量。Hitchin 的构造取 $M=T^*\mathcal{A}$（紧 Riemann 面上联络空间 $\mathcal{A}$ 的余切丛）、$G$ 为规范群，此时 $\mu=0$ 恰是 Hitchin 方程，商得到 $M_H$（稳定部分）。有限维玩具模型是 $T^*G$ 与余伴随轨道。

\textbf{在 K--W 中的位置}：§4 用超凯勒商刻画 $M_H$（N4、S3），把 4d 方程的模空间与 Hitchin 方程的解空间等同起来；这一步是全文从“规范理论”过渡到“模空间几何”的铰链。

\textbf{读原文前需要的前置}：注释 3。

\textbf{常见坑}：这里的商是无限维的，需要 Sobolev 完备化与 slice 定理，不能当有限维 GIT 商直接操作；商之前必须固定稳定条件，否则商空间非 Hausdorff（轨道可能聚在一起）。

\noindent\textbf{链接}：\url{https://ncatlab.org/nlab/show/symplectic+reduction}
\end{kcard}

\begin{kcard}{注释 5：Higgs 丛（Higgs bundle）}
\anchor{nLab: Higgs+bundle}

\textbf{一句话}：Higgs 丛是全纯向量丛 $E$ 加上一个取值于 $K_X$ 的 Higgs 场 $\Phi$，并满足 $\Phi\wedge\Phi=0$。

\textbf{数学要点}：Higgs 丛是二元组 $(E,\Phi)$，其中 $E\to X$ 是全纯向量丛，$\Phi\in H^0(X,\mathrm{End}(E)\otimes\Omega^1_X)$ 是 $\mathcal{O}_X$-线性的“扭化自同态”，并满足可积条件 $\Phi\wedge\Phi=0$（$X$ 为曲线时自动成立）。稳定性只对 $\Phi$-不变子丛检验斜率，这与普通丛稳定性不同。$\deg=0$ 的稳定 Higgs 丛模空间就是 $M_H$。秩 $1$ 时 $\Phi$ 是任意全纯 $1$-形式，$(E,\Phi)$ 对应余切丛 $T^*\mathrm{Jac}(X)$ 的点，可作为最小模型。与平坦联络的差别正在于 $\Phi$ 无 $\dd\Phi$ 项：它是把 $\bar\partial$ 与 $\Phi$ 拼成复化联络后剩下的“额外一阶项”。

\textbf{在 K--W 中的位置}：§4 给出 $M_H$ 的 Dolbeault（Higgs 丛）侧（N4、S3）；§8 的电特征膜用局部系统/平坦丛描述，与这里的 Higgs 侧构成非阿贝尔 Hodge 的两面（N8）；§11 的 A-膜与 D-模也以 $M_H$ 上的几何为舞台（N11）。

\textbf{读原文前需要的前置}：注释 1、注释 3。

\textbf{常见坑}：$\Phi\wedge\Phi=0$ 只在 $\dim X>1$ 时非平凡；稳定性必须只检验 $\Phi$-不变子丛，否则定义过强、几乎无稳定对象；B-模型的对象不是“任意全纯丛”，必须带 $\Phi$。

\noindent\textbf{链接}：\url{https://ncatlab.org/nlab/show/Higgs+bundle}
\end{kcard}

\begin{kcard}{注释 6：Hitchin 模空间（Hitchin moduli space）}
\anchor{nLab: Higgs+bundle}

\textbf{一句话}：$M_H$ 是稳定 Higgs 丛的模空间，也是 K--W 全文的舞台：超凯勒结构、可积系统、几何 Langlands 都发生在它身上。

\textbf{数学要点}：$M_H$ 的复维数是 $\dim_{\mathbb{C}}M_H=2\dim_{\mathbb{C}}\mathrm{Bun}_G=2(g-1)\dim G$。作为实流形，它同时是 de Rham 侧（平坦丛/局部系统）与 Dolbeault 侧（Higgs 丛）的模空间，两套复结构由非阿贝尔 Hodge 对应联系起来。三个复结构 $I,J,K$ 中，$J$ 下 $M_H$ 是 Hitchin 可积系统（纤维化），$I$ 与 $K$ 则用于不同的 B-模型。$M_H$ 又是 4d $N=4$ 超杨--米尔斯紧化到 2d 后 sigma 模型的靶空间，其超凯勒度量给出拓扑场论的作用量。$N=4$ 的 S-对偶在 $M_H$ 上交换复结构与膜类型，这正是几何 Langlands 的物理来源。

\textbf{在 K--W 中的位置}：§4 紧化得到 $M_H$（N4、S3）；§5 的 2d 拓扑场论以 $M_H$ 为靶空间（N5）；§7--§12 的所有膜都是 $M_H$ 的子流形或 $M_H$ 上的对象（N7--N12）。

\textbf{读原文前需要的前置}：注释 3、注释 5。

\textbf{常见坑}：“$M_H$ 上全纯”必须说明是哪个复结构（$I$、$J$ 还是 $K$），否则句子没有意义；物理紧化里的 $M_H$ 与数学上紧超凯勒商得到的空间在紧化/稳定性上有差别，不要混为一谈；不同群 $G$ 的 $M_H$ 维数与奇异结构不同。本文的约定是：复结构 $I$ 下 $M_H$ 是 Higgs 丛模空间（Hitchin 可积系统在这里），复结构 $J$、$K$ 下它是局部系统/平坦丛的模空间。

\noindent\textbf{链接}：\url{https://ncatlab.org/nlab/show/Higgs+bundle}
\end{kcard}

\begin{kcard}{注释 7：Hitchin 方程（Hitchin equations）}
\anchor{nLab: Higgs+bundle}

\textbf{一句话}：Hitchin 方程是联络与 Higgs 场的一阶自对偶方程，其解模去规范群正是 $M_H$。

\textbf{数学要点}：在紧 Riemann 面 $X$ 上取 Hermitian 丛 $(E,\bar\partial_A)$ 与 $\Phi\in\Omega^{1,0}(\mathrm{End}\,E)$，Hitchin 方程是
\[
F_A+[\Phi,\Phi^*]=0,\qquad \bar\partial_A\Phi=0 .
\]
这是 4d 自对偶 Yang--Mills 方程沿 Kähler 方向降维的产物；其解空间模规范群同构于稳定 Higgs 丛模空间（Hitchin 1987，Simpson 1992）。用超凯勒语言，它就是矩映射 $\mu=0$ 的三个分量：$\mu_{\mathbb{R}}=F_A+[\Phi,\Phi^*]$ 与 $\mu_{\mathbb{C}}=\bar\partial_A\Phi$。它与 §10 的 Bogomolny 方程同属“降维自对偶”家族，只是一个在曲面上、一个在 3 维空间中。

\textbf{在 K--W 中的位置}：§4（N4、S3）把 4d 方程紧化到 2d 后得到它，$M_H$ 就是其解模空间；§10 的 Bogomolny 方程作为单极子侧类比（N10）。后文使用“解 Hitchin 方程”这一说法时，指的就是这里的方程组。

\textbf{读原文前需要的前置}：注释 3、注释 4、注释 5。

\textbf{常见坑}：两个方程缺一不可——只写 $F_A+[\Phi,\Phi^*]=0$ 并不强制 $\Phi$ 全纯；$\Phi^*$ 依赖 Hermitian 度量，方程规范不变但解依赖度量；解的存在性等价于稳定性（Kobayashi--Hitchin 型定理），并非每个丛都有解。

\noindent\textbf{链接}：\url{https://ncatlab.org/nlab/show/Higgs+bundle}
\end{kcard}

\begin{kcard}{注释 8：Hitchin 纤维化与可积系统（Hitchin fibration and integrable systems）}
\anchor{nLab: Hitchin+fibration}

\textbf{一句话}：Hitchin 纤维化把 $M_H$ 送到“特征多项式”空间，使它成为 Liouville 可积系统，一般纤维是谱曲线的 Jacobi/Prym 簇。

\textbf{数学要点}：对 $(E,\Phi)$ 取谱曲线 $\det(\lambda-\Phi)=0$，它给出分支覆盖 $\pi:S\to X$；特征多项式的系数定义 Hitchin 基 $\mathcal{B}=\bigoplus_{i=2}^{r}H^0(X,K_X^{\otimes i})$（$\mathrm{SL}(r)$ 情形没有 $i=1$ 项）。 Hitchin 映射 $h:M_H\to\mathcal{B}$ 是 proper 的、Lagrangian 的代数完全可积系统；一般纤维是谱曲线的 Prym 簇，是阿贝尔簇，维数恰为 $\dim M_H/2$。$\mathcal{B}$ 上的坐标给出两两泊松对易的 Hamiltonians，因此是可积系统。Beilinson--Drinfeld 用它做量子化并连到几何 Langlands，Ngô 则用 Hitchin 纤维化证明 fundamental lemma。

\textbf{在 K--W 中的位置}：§4 与 §9 用谱曲线、Hitchin 基讨论 $M_H$ 的结构（N4、N9）；$M_H$ 在复结构 $J$ 下的可积系统是后文讨论谱曲线、BPS 态与单极子时的数学原型（§10，N10）。

\textbf{读原文前需要的前置}：注释 3、注释 6。

\textbf{常见坑}：只有一般纤维才是阿贝尔簇，判别式轨迹上的奇异纤维携带关键信息；不同群（$\mathrm{SL}$ 与 $\mathrm{PGL}$）的 Hitchin 基不同，Langlands 对偶会交换这两者；“可积”是 Liouville 意义下的，不是量子可积。

\noindent\textbf{链接}：\url{https://ncatlab.org/nlab/show/Hitchin+fibration}
\end{kcard}

\begin{kcard}{注释 9：非阿贝尔 Hodge 对应（nonabelian Hodge theory）}
\anchor{nLab: nonabelian+Hodge+theory}

\textbf{一句话}：非阿贝尔 Hodge 对应断言：紧 Kähler 流形上“平坦丛”与“稳定 Higgs 丛”是同一个模空间的两套复结构。

\textbf{数学要点}：阿贝尔情形是 Dolbeault 分解 $H^1(X,\mathbb{C})\cong H^1(X,\mathcal{O}_X)\oplus H^0(X,\Omega^1_X)$。非阿贝尔推广说：$c_1=0$ 的稳定 Higgs 丛与不可约平坦联络一一对应（Simpson 1992；Corlette 的调和度量定理是分析核心）。两侧模空间 $M_{\mathrm{dR}}$ 与 $M_{\mathrm{Dol}}$ 实解析同胚、复结构不同，正如 $M_H$ 上的 $I$ 与 $J$。用栈的语言即 $H(X_{\mathrm{dR}},BG)\cong H(X_{\mathrm{Dol}},BG)^{ss,0}$。这个对应解释了 $M_H$ 为何同时是平坦丛空间与 Higgs 丛空间，也是几何 Langlands 的经典极限：S-对偶在其上交换 de Rham 侧与 Dolbeault 侧。

\textbf{在 K--W 中的位置}：§4 与 §8--§11（N4、N8、N11）；§8 的电特征膜把 Dolbeault 侧的 B-膜对应到 de Rham 侧的平坦丛，正是这个对应的物理化身；骨架 S6 与 S8（局部系统 $\leftrightarrow$ D-模）。

\textbf{读原文前需要的前置}：注释 5、注释 6、注释 11。

\textbf{常见坑}：必须限制在 $c_1=0$（$\deg=0$）且不可约/稳定的部分；两侧同胚是实解析的，不是复解析同构；非紧 Kähler 情形需要额外的解析假设。

\noindent\textbf{链接}：\url{https://ncatlab.org/nlab/show/nonabelian+Hodge+theory}
\end{kcard}

\begin{kcard}{注释 10：广义复几何与广义复结构（generalized complex geometry and generalized complex structure）}
\anchor{nLab: generalized+complex+geometry}

\textbf{一句话}：广义复结构是把复结构与辛结构装进同一个 $J$ 里：它在 $TX\oplus T^*X$ 上定义，平方为 $-1$，且关于自然配对反对称。

\textbf{数学要点}：广义切丛是 $TX\oplus T^*X$，上面有自然配对 $\langle (v,\xi),(w,\eta)\rangle=\xi(w)+\eta(v)$。广义复结构是 $\mathcal{J}:TX\oplus T^*X\to TX\oplus T^*X$，满足 $\mathcal{J}^2=-1$ 且 $\mathcal{J}^*=-\mathcal{J}$（关于上述配对）；后一条使它同时是复结构与辛结构。两个基本例子是 $\mathcal{J}_J=\mathrm{diag}(-J,J^*)$（来自复结构 $J$）与 $\mathcal{J}_\omega=\begin{pmatrix}0&-\omega^{-1}\\ \omega&0\end{pmatrix}$（来自辛形式 $\omega$）。逐点来看，$\mathcal{J}$ 给出结构群的约化 $U(n,n)\hookrightarrow O(2n,2n)$（Gualtieri 命题 4.6）。这套几何由 Hitchin 提出、Gualtieri 与 Cavalcanti 发展，统一了复几何、辛几何与黎曼几何，其物理动机是 $T$-对偶与超引力；它与镜像对称、Poisson/Courant sigma 模型都相关。

\textbf{在 K--W 中的位置}：§5 的 2d 拓扑场论从广义复几何的语言进入（N5）；$M_H$ 上不同复结构对应不同广义复结构，从而决定 twist 与膜的类型。§11--§12 讨论膜与 D-模时，广义复几何提供“B-膜 = 复子流形、A-膜 = 拉格朗日子流形”的统一框架。

\textbf{读原文前需要的前置}：注释 3、注释 5。

\textbf{常见坑}：$\mathcal{J}$ 不是 $TX$ 上的复结构，作用对象是 $TX\oplus T^*X$；$\mathcal{J}^2=-1$ 单独不够，还必须要求它对自然配对反对称，否则不给出辛结构；广义复流形不一定是复流形，反之亦然。

\noindent\textbf{链接}：\url{https://ncatlab.org/nlab/show/generalized+complex+geometry}
\end{kcard}

\begin{kcard}{注释 11：平坦联络与局部系统（flat connection and local system）}
\anchor{nLab: flat+connection}

\textbf{一句话}：平坦联络就是曲率为零的联络；沿着它平行移动只依赖同伦类，于是丛的“数据”退化成基本群的一个表示，这就是局部系统。

\textbf{数学要点}：纤维丛上的联络 $A$ 称为平坦的，若曲率 $F=\dd A+A\wedge A=0$；等价说法是 Maurer--Cartan 方程成立，在可积系统里也叫零曲率方程（Lax 方程总能写成这个形式）。平坦联络的平行移动与路径的同伦类无关，因此给出从基本群胚到结构群的表示，也就是局部系统；在这个意义下“平坦联络 = 局部系统 = 基本群的表示”就是 Riemann--Hilbert 对应的最朴素的形态。复流形上有 Koszul--Malgrange 定理：全纯平坦联络与全纯向量丛（带平坦联络）一致。紧 Riemann 面上，向量丛容许平坦联络当且仅当 $|\langle c_1(E),[X]\rangle|<g$（Milnor）；$\deg=0$ 时平坦联络模空间与稳定丛模空间的对应即 Narasimhan--Seshadri 定理。

\textbf{在 K--W 中的位置}：§8 的电特征膜（electric eigenbrane）在 de Rham 侧就是带平坦联络的丛或局部系统（N8、S6）；§4 的 $M_H$ 在复结构 $J$、$K$ 下是局部系统模空间，正是靠这个等同。Hecke 特征层在 de Rham 侧被表述为“Hecke 算子作用在局部系统上取本征”。

\textbf{读原文前需要的前置}：注释 1、注释 3。

\textbf{常见坑}：平坦是曲率为零，不是“联络取零”；平凡联络平坦，但平坦联络一般不平凡；局部系统是局部常值层，不是“局部定义的联络”；正特征或非紧情形需要额外假设。

\noindent\textbf{链接}：\url{https://ncatlab.org/nlab/show/flat+connection}
\end{kcard}

\begin{kcard}{注释 12：Riemann--Hilbert 对应（Riemann--Hilbert correspondence）}
\anchor{nLab: Riemann-Hilbert+correspondence}

\textbf{一句话}：Riemann--Hilbert 对应说，线性微分方程的解系与拓扑数据是同一件事：正则完整 D-模的导出范畴等价于构造可层（constructible sheaves）的导出范畴。

\textbf{数学要点}：基本形式是“平坦向量丛的范畴 $\simeq$ 局部系统的范畴”。在代数几何里推广为 de Rham 函子 $\mathrm{dR}:D(X_{\mathrm{Zar}},\Dmod_X)\to D(X_{\mathrm{an}},\mathbb{C})$，把 D-模送到其解复形（$\Omega^\bullet(\mathcal{E}^{\mathrm{hol}})$）。Mebkhout、Kashiwara 与 Beilinson--Bernstein 的定理说：$\mathrm{dR}$ 限制在正则完整 D-模的导出范畴上给出等价 $D^b_{\mathrm{rh}}(X_{\mathrm{Zar}},\Dmod_X)\xrightarrow{\sim}D^b_c(X_{\mathrm{an}},\mathbb{C})$。非正则情形要换成 ind-sheaves（D'Agnolo--Kashiwara）。在几何 Langlands 中，这个对应把自守侧从“反常层/构造可层”翻译成“D-模”，从而能与 Hecke 算子、谱侧对偶起来（Frenkel 2005，§3.4--3.6）。

\textbf{在 K--W 中的位置}：§11 的 A-膜与 D-模（N11、S8）：A-膜范畴被期待等价于 D-模范畴，Riemann--Hilbert 是这一步的拓扑侧；§12 把 Hecke 特征层实现为 D-模，也依赖它把“局部系统”与“D-模”放在同一张桌上。

\textbf{读原文前需要的前置}：注释 11、注释 28。

\textbf{常见坑}：等价只在正则完整（regular holonomic）子范畴上成立，并非所有 D-模；$\mathrm{dR}$ 不是“取解空间”这种朴素函子，而是导出函子；复解析化 $X_{\mathrm{an}}$ 不能省。

\noindent\textbf{链接}：\url{https://ncatlab.org/nlab/show/Riemann-Hilbert+correspondence}
\end{kcard}

\begin{kcard}{注释 13：A-模型与 B-模型（A-model and B-model）}
\anchor{nLab: A-model}

\textbf{一句话}：A-模型与 B-模型是同一个 2d 超对称 sigma 模型的两个拓扑 twist：A 只看辛结构（Fukaya 范畴、Gromov--Witten 理论），B 只看复结构（相干层的导出范畴）。

\textbf{数学要点}：A-模型是 2 维拓扑共形场论（TCFT），其 Calabi--Yau 范畴是辛流形 $(X,\omega)$ 的 Fukaya 范畴；它等于底层 Poisson sigma 模型做规范固定后的 AKSZ 理论（AKSZ 1997），有效上就是 $X$ 的 Gromov--Witten 理论。物理来源是超弦在 Calabi--Yau 上传播：对 2d 超对称 sigma 模型的顶点算子代数作形变，使一个超 Virasoro 生成元的平方为零，得到的“拓扑 twist”代数正是 TCFT 的 BRST 复形。B-模型则是复结构侧的拓扑理论，其范畴是相干层的导出范畴 $D^b_{\infty}(X)$。两者由同调镜像对称联系起来：$\mathrm{Fuk}(X)\simeq D^b_\infty(\widehat X)$。A-模型的开弦 D-膜由 Fukaya 范畴给出。

\textbf{在 K--W 中的位置}：§5 的 2d 拓扑场论就是这两种模型（N5）；在 K--W 中把 4d 理论紧化在 $C$ 上，低能有效理论在特定参数下正是靶空间为 $M_H$ 的 A-模型与 B-模型（N4、N11）。§11 之后所有“A-膜/B-膜”的说法都相对于这两个模型而言。

\textbf{读原文前需要的前置}：注释 10、注释 14。

\textbf{常见坑}：A/B 不是两个不同的物理理论，而是同一理论的两种 twist；A-模型依赖辛结构而非复结构，B-模型相反；说“B-模型的对象是相干层”时要指明是导出范畴而不是普通范畴。

\noindent\textbf{链接}：\url{https://ncatlab.org/nlab/show/A-model}；\url{https://ncatlab.org/nlab/show/B-model}
\end{kcard}

\subsection*{第二组\quad 物理语言：膜、扭化、对偶与 BPS}

\begin{kcard}{注释 14：镜像对称与同调镜像对称（mirror symmetry and homological mirror symmetry）}
\anchor{nLab: mirror+symmetry}

\textbf{一句话}：镜像对称说两个 Calabi--Yau 流形 $X$ 与 $\widehat X$ 的 A-模型与 B-模型互换；同调镜像对称把它精确化为两个 Calabi--Yau 范畴的等价。

\textbf{数学要点}：对复 3 维 Calabi--Yau 流形 $X$，存在（至少对一部分 $X$）镜像 $\widehat X$，交换 Hodge 数 $h^{1,1}\leftrightarrow h^{1,2}$，使得 $\mathrm{SCFT}_A(X)\simeq\mathrm{SCFT}_B(\widehat X)$；在拓扑层面即 $A(X)\simeq B(\widehat X)$、$B(X)\simeq A(\widehat X)$。Kontsevich 1994 给出范畴化表述：Fukaya 范畴与相干层导出范畴等价，$\mathrm{Fuk}(X)\simeq D^b_\infty(\widehat X)$，这就是同调镜像对称（HMS）；Kontsevich 与 Costello 进一步证明 TCFT 的数据等价于 Calabi--Yau $A_\infty$-范畴的数据，使两种表述一致。镜像对称可理解为 Strominger--Yau--Zaslow 意义下 Lagrangian 纤维化上的 $T$-对偶。

\textbf{在 K--W 中的位置}：§5--§6 讨论 2d 拓扑场论与 S-对偶时，镜像对称是核心机制（N5、N6、S4）：K--W 的 S-对偶在 $M_H(G)$ 与 $M_H(\Ldual)$ 之间表现为镜像对称，交换 B-膜与 A-膜。§12 的几何 Langlands 对应因此可以读作“$M_H$ 上的镜像对称”。

\textbf{读原文前需要的前置}：注释 13、注释 10。

\textbf{常见坑}：镜像对称不是“两个流形同构”，而是两个理论等价；HMS 是范畴等价而非流形层面的对应；镜像的构造对一般 Calabi--Yau 仍是猜想，不要当成已证定理。

\noindent\textbf{链接}：\url{https://ncatlab.org/nlab/show/mirror+symmetry}
\end{kcard}

\begin{kcard}{注释 15：D-膜与 A/B-膜（D-branes and A/B-branes）}
\anchor{nLab: A-brane}

\textbf{一句话}：D-膜是开弦的边界条件；在拓扑 twist 之后，A-模型里的膜是辛几何对象（Lagrangian/余各向同性地），B-模型里的膜是全纯几何对象（复子流形、相干层）。

\textbf{数学要点}：D-膜是开弦端点必须停留的子流形，其世界体积上还带有 Chan--Paton 丛与联络；从 CFT 角度，它是一个保持边界条件的边界态。把 2d 超对称 sigma 模型作拓扑 twist 后，D-膜成为拓扑膜：A-模型的 D-膜（A-膜）构成 Fukaya 范畴，其典型支撑是 Lagrangian 子流形（更一般地是余各向同性子流形，见注释 16），对象之间的态射由 Floer 上同调给出；B-模型的 D-膜（B-膜）则是复子流形上的相干层。同调镜像对称把 A-膜范畴与 B-膜范畴互换。

\textbf{在 K--W 中的位置}：§7 的线算子被实现为膜（N7、N12）；§8 的电特征膜与磁特征膜分别是 B-膜与 A-膜（N8）；§9 的 't Hooft 算子被实现为 Hecke 算子，也通过膜的语言表述（N9、S7）。全文“三种膜”的框架（N12）就是 A-膜、B-膜与它们的特征对象。

\textbf{读原文前需要的前置}：注释 13、注释 14。

\textbf{常见坑}：D-膜不是“弦的端点”，而是端点可停留的整个子流形及其上的丛；B-膜一般不是复子流形而是相干层（可能带扭）；A-膜也不一定支撑在 Lagrangian 上，余各向同性支撑同样合法。

\noindent\textbf{链接}：\url{https://ncatlab.org/nlab/show/A-brane}；\url{https://ncatlab.org/nlab/show/B-brane}
\end{kcard}

\begin{kcard}{注释 16：余各向同性子流形与典范余各向同性膜（coisotropic submanifold and canonical coisotropic brane）}
\anchor{nLab: coisotropic+submanifold}

\textbf{一句话}：余各向同性子流形是“反对称张量场把它推到切空间里”的子流形；在 K--W 中，覆盖整个 $M_H$ 的典范余各向同性膜是把 A-膜变成 D-模的那台机器。

\textbf{数学要点}：设 $(X,\pi)$ 是 Poisson 流形，子流形 $S\hookrightarrow X$ 称为余各向同性的，若用 Poisson 张量收缩余法丛的映射 $\pi:T^*X\to TX$ 限制在 $N^*S$ 上穿过 $TS$：$\pi(N^*S)\subset TS$；等价地，在 $S$ 上 vanish 的函数构成的子代数对 Poisson 括号封闭。辛流形是 Poisson 流形的特例，此时 $S$ 余各向同性的条件是 $TS$ 的辛正交补 $TS^{\perp}\subset TS$（即 $\dim S\ge\dim X/2$，最极端情形是 Lagrangian）。A-膜可以支撑在余各向同性子流形上，条件是 Chan--Paton 曲率 $F$ 满足 $(\omega^{-1}F)^2=-1$ 在 $N^*S$ 上成立（Kapustin--Orlov）。K--W 的典范余各向同性膜 $\mathcal{A}_{cc}$ 支撑在整个 $M_H$ 上，其曲率就是 Kähler 形式，用它作“探针”可以把 A-膜范畴与 D-模范畴联系起来。

\textbf{在 K--W 中的位置}：§9 与 §12 用 $\mathcal{A}_{cc}$ 上的 $\mathrm{Hom}(\mathcal{A}_{cc},\mathcal{A})$ 构造 D-模，并预期得到 Hecke 特征层（N9、N12、S7）；§11 的“A-膜 = D-模”对应也靠它实现（N11）。

\textbf{读原文前需要的前置}：注释 10、注释 15。

\textbf{常见坑}：余各向同性包含 Lagrangian 作为特例，但比 Lagrangian 大；条件用的是余法丛而非法丛；$\mathcal{A}_{cc}$ 支撑在整个空间上，不是“子流形上的膜”；这里的曲率条件是 $(\omega^{-1}F)^2=-1$，别与 Hermitian--Yang--Mills 方程混淆。

\noindent\textbf{链接}：\url{https://ncatlab.org/nlab/show/coisotropic+submanifold}；\url{https://ncatlab.org/nlab/show/A-brane}
\end{kcard}

\begin{kcard}{注释 17：拓扑扭化（topological twist）}
\anchor{nLab: topological+twist}

\textbf{一句话}：拓扑扭化是把超对称理论的 R-对称群与旋转群混合，挑出一个标量超荷 $Q$，再把理论换成它的 $Q$-上同调——于是理论变成拓扑场论。

\textbf{数学要点}：一般程序是：超对称理论里的超荷按某种表示变换；把 R-对称的一个子群嵌入到旋转群中，使某个超荷变成标量且满足 $Q^2=0$；然后把 Hilbert 空间换成 $Q$ 的上同调，关联函数就只依赖时空的拓扑而非度规。Witten 1988 用这个办法从 2d 超对称 sigma 模型造出 A-模型与 B-模型；对 4d $N=4$ 超杨--米尔斯，Witten 1991 的扭化把 $SU(4)_R$ 与 $SO(4)$ 旋转群组合，得到一个标量拓扑超荷。K--W 的扭化还带一个 $\mathbb{CP}^1$ 参数 $t$：超荷取两个分量的组合，$(u:v)\in\mathbb{CP}^1$ 标记这一族拓扑理论。\anchor{src p.16}

\textbf{在 K--W 中的位置}：§3 的核心操作（N3、S2、M1）：扭化后的拓扑超荷 $Q$ 满足 $Q^2=0$，它把 $N=4$ 超杨--米尔斯变成一个 4d 拓扑场论；§5 在 $C$ 上紧化后得到的 2d 理论又是对同一扭化的进一步处理（N5、S4）。五个 miracle 里的 M1 说的正是：整族 $\mathbb{CP}^1$ 参数下拓扑性都保持，只有 $t=\pm\ii$ 处退化。

\textbf{读原文前需要的前置}：注释 18、注释 20。

\textbf{常见坑}：扭化不改作用量，改的是“旋转群怎么嵌进超对称代数”，所以不要把它误当成加一个拓扑项；不是任意超荷都能用来扭化，必须是标量且平方为零的组合；$t=\pm\ii$ 的退化点不是“额外假设”，而是 K--W 论证的关键结构。

\noindent\textbf{链接}：\url{https://ncatlab.org/nlab/show/topological+twist}
\end{kcard}

\begin{kcard}{注释 18：$N=4$ 超杨--米尔斯（$N=4$ super Yang--Mills theory）}
\anchor{nLab: D4 N4 super Yang-Mills theory}

\textbf{一句话}：4 维 $N=4$ 超杨--米尔斯是四维最大的超对称规范理论，整个多重态只有规范场、4 个 Weyl 费米子与 6 个标量，全部取伴随表示，并且有 $SL(2,\ZZ)$ 电磁对偶。

\textbf{数学要点}：$N=4$ 指 4 组超荷，共 16 个超荷；R-对称群是 $SU(4)\cong \mathrm{Spin}(6)$，恰好是 6 个标量所在的空间的旋转群。理论共形（$\beta$ 函数为零），且是唯一同时具有最大超对称与共形不变性的四维规范理论；它可看作 6d $(2,0)$ 超共形场论在环面上的紧化。S-对偶（Montonen--Olive）作用在复耦合 $\tau=\theta/2\pi+4\pi\ii/g^2$ 上：$\tau\mapsto-1/\tau$，同时交换电与磁、把规范群 $G$ 换成 Langlands 对偶群 $\Ldual$；$\tau\mapsto\tau+1$ 是 $\theta$ 的平移。合起来给出 $SL(2,\ZZ)$ 作用。

\textbf{在 K--W 中的位置}：§2 的出发点（N2、S1）：K--W 的第一句话就是“物理的电磁对偶群 $SL(2,\ZZ)$ 正是 Langlands 对偶的那个群”；§3 的扭化以它为原料（N3）；§4 把它放在 $\Sigma\times C$ 上紧化得到 $M_H$ 上的 sigma 模型（N4）。

\textbf{读原文前需要的前置}：注释 1、注释 17。

\textbf{常见坑}：$N=4$ 的“4”是超荷组数，不是场的数目（场有 $1+4+6$ 个分量，都在伴随表示）；S-对偶不是把 $g$ 换成 $1/g$ 的微扰展开，而是非微扰对偶；$SL(2,\ZZ)$ 作用在 $\tau$ 上，与 Langlands 对偶群 $\Ldual$ 是两个不同角色的群，不要混为一谈。

\noindent\textbf{链接}：\url{https://ncatlab.org/nlab/show/D%3D4+N%3D4+super+Yang-Mills+theory}
\end{kcard}

\begin{kcard}{注释 19：BPS 态（BPS state）}
\anchor{nLab: BPS+state}

\textbf{一句话}：BPS 态是被一部分超荷湮灭、从而质量恰好等于中心荷（饱和 BPS 界）的态；它的多重态比一般态短，因此谱受保护、可以从弱耦合外推。

\textbf{数学要点}：在扩展超对称代数中，超荷的反交换子形如 $\{Q_A,Q_B\}=E\,\delta_{AB}-K_{AB}$，其中 $E$ 是能量/质量、$K_{AB}$ 是中心荷；右边正定性给出 $E\ge |K|$，这就是 BPS 界。当某些本征值取等号时，对应的超荷组合湮灭该态，态落在“短多重态”里，称为 BPS 态（Bogomolny--Prasad--Sommerfield 饱和）。由于超对称保护，BPS 态的数量与耦合常数无关，这是把弱耦合计数外推到强耦合的核心工具；4d 磁单极子、瞬子、以及 4d $N=2$ 的谱网络都与 BPS 谱有关，Donaldson--Thomas 与 Gromov--Witten 不变量也可看成 BPS 态的计数。

\textbf{在 K--W 中的位置}：§10 讨论 Bogomolny 方程与 Hecke 修改空间时，单极子冒泡与 BPS 谱是物理图像（N10、M3）；§11 的 't Hooft 算子本征值也可以读作 BPS 谱的代数化。

\textbf{读原文前需要的前置}：注释 23、注释 18。

\textbf{常见坑}：BPS 是“饱和超对称界”，不是“稳定”；$1/2$-BPS 指保留一半超荷，比例要先数清楚再谈；没有中心荷时 BPS 界退化为平凡不等式（如纯 $N=4$ 的无荷扇区），此时不要滥用“BPS 态”。

\noindent\textbf{链接}：\url{https://ncatlab.org/nlab/show/BPS+state}
\end{kcard}

\begin{kcard}{注释 20：BRST 上同调（BRST cohomology）}
\anchor{nLab: BRST+complex}

\textbf{一句话}：BRST 上同调是用一个幂零算子 $Q$ 取规范理论物理态的标准机器；在拓扑扭化之后，这台机器取的是拓扑不变量。

\textbf{数学要点}：量子化规范理论时引入 ghost 场 $c$（次数 1），构造 BRST 算子 $Q_{\mathrm{BRST}}$，满足 $Q_{\mathrm{BRST}}^2=0$；物理态定义为 $H^0(Q_{\mathrm{BRST}})$。从数学上看，BRST 复形是模空间（或相应 $\infty$-groupoid）的 Chevalley--Eilenberg 代数：ghost 对应规范变换，更高次的 ghost-of-ghost 对应高阶规范变换；BV 形式是它在场论上的完整推广。拓扑扭化后，原有的某个超荷变成标量且平方为零，把 Hilbert 空间换成它的上同调，关联函数就变成拓扑不变量（Witten 1988 的 A/B 模型、Witten 1991 的 4d 拓扑杨--米尔斯都是这个套路）。K--W 的 4d TQFT 正是由扭化超荷 $Q$ 定义的，Wilson 算子、't Hooft 算子与各种膜都是 $Q$-上同调类。

\textbf{在 K--W 中的位置}：§3 的拓扑超荷 $Q$ 就是全文的“上同调算子”（N3、S2）；§4--§5 紧化后的 2d 理论继承同一个 $Q$（N4、N5），因此 2d 侧的膜范畴是 $Q$-上同调意义下的范畴。

\textbf{读原文前需要的前置}：注释 17、注释 18。

\textbf{常见坑}：BRST 上同调不等于“把规范固定掉”那么简单，全局与非微扰信息仍在；扭化得到的 $Q$ 来自超对称，与通常意义下的 BRST 算子来源不同，不要混用记号；取 $H^0$ 还是整个上同调取决于理论是 cohomological 还是 quantum 的。

\noindent\textbf{链接}：\url{https://ncatlab.org/nlab/show/BRST+complex}
\end{kcard}

\begin{kcard}{注释 21：瞬子与 ADHM 构造（instantons and the ADHM construction）}
\anchor{nLab: ADHM+construction}

\textbf{一句话}：Yang--Mills 瞬子是自对偶联络，ADHM 构造把 $\mathbb{R}^4$（或 $S^4$）上给定拓扑荷的全部瞬子写成有限维线性代数数据再作超凯勒商，于是瞬子模空间变成一个可以计算的对象。

\textbf{数学要点}：瞬子解满足自对偶方程 $F=*F$，它使 Yang--Mills 作用量在同伦类中取极小，拓扑荷为 $c_2=k$。ADHM 数据取复向量空间 $V\cong\CC^k$、$W\cong\CC^N$ 以及 $A,B\in\operatorname{End}(V)$、$I\in\operatorname{Hom}(W,V)$、$J\in\operatorname{Hom}(V,W)$，满足方程 $[A,B]+IJ=0$ 与 $[A,A^\dagger]+[B,B^\dagger]+II^\dagger-J^\dagger J=0$，再按 $U(k)$ 的作用作商。这个商恰好是超凯勒商（Kronheimer 的表述）：两个方程分别是复与实矩映射的零值，因此瞬子模空间继承超凯勒结构。ADHM 把瞬子问题化归为矩阵方程，从而给出模空间的维数、指标定理计数与 Donaldson 不变量计算的基础；其原始文献是 Atiyah--Hitchin--Drinfeld--Manin（1978），同期的 Christ--Weinberg--Stanton 给出独立构造。

\textbf{在 K--W 中的位置}：§3--§4 的 4d $N=4$ 理论里，瞬子是规范场的基本非微扰对象（N3、N4），拓扑扭化后的路径积分正是按瞬子荷展开的；§9--§10 的单极子冒泡与 Hecke 修改空间在结构上平行于瞬子模空间（N9、N10）。

\textbf{读原文前需要的前置}：注释 3、注释 4。

\textbf{常见坑}：ADHM 构造针对 $\mathbb{R}^4$ 或 $S^4$（平坦或紧致化）上的瞬子，一般流形要用别的构造；那个 $U(k)$ 商是超凯勒商而不是朴素的 GIT 商，别按稳定丛的直觉去处理；拓扑荷 $k$ 与规范群的中心、框架化（framing）约定有关，比较不同文献的模空间维数前先核对约定。

\noindent\textbf{链接}：\url{https://ncatlab.org/nlab/show/ADHM+construction}
\end{kcard}

\begin{kcard}{注释 22：磁单极子与 Nahm 方程（magnetic monopoles and Nahm's equation）}
\anchor{nLab: Nahm's equation}

\textbf{一句话}：'t~Hooft--Polyakov 磁单极子是 3d 规范理论中带拓扑荷的孤子解，而 Nahm 方程是一组矩阵值常微分方程，它通过 Nahm 变换把单极子模空间上的数据对偶地编码成另一组线性数据。

\textbf{数学要点}：在带伴随标量 $\phi$ 的规范理论中，Bogomolny 方程 $F=*D\phi$ 的单极子解具有质量 $4\pi v|n|$、磁荷 $n$；固定 $n$ 的模空间是超凯勒流形（Hitchin 的度量）。Nahm 数据是三个 $\mathfrak{u}(N)$-值函数 $X^1,X^2,X^3$（$y\in(0,\infty]$），满足
\[
\frac{dX^3}{dy}+[X^1,X^2]=0,\qquad
\frac{dX^1}{dy}+[X^2,X^3]=0,\qquad
\frac{dX^2}{dy}+[X^3,X^1]=0,
\]
在 $y\to0$ 处有极点解 $X^i(y)=\rho^i/y+\text{非奇异项}$，其中 $\rho:\mathfrak{su}(2)\to\mathfrak{u}(N)$。Nahm 变换（Atiyah--Hitchin--Drinfeld--Manin--Nahm）在单极子与 Nahm 数据之间建立对偶；Donaldson（1984）用它分类了单极子。极点条件在现代文献中称为 Nahm 极点边界条件（Mazzeo--Witten 2014），它是超对称规范理论中一类重要边界条件。

\textbf{在 K--W 中的位置}：§9--§10 讨论 't~Hooft 算子与 Hecke 修改时，单极子冒泡（monopole bubbling）是核心图像，边界附近的 BPS 方程退化为 Nahm 型方程（N9、N10、M3）；§12 的边界条件构造中也会再次遇到 Nahm 极点（N12）。

\textbf{读原文前需要的前置}：注释 18、注释 19。

\textbf{常见坑}：Nahm 方程不是 Yang--Mills 方程本身，而是降维/边界条件方程；Nahm 变换要在对偶环面上做，别把两边的规范群与拓扑荷记混；Nahm 极点边界条件与普通 Dirichlet 边界条件不同，它允许场在边界附近发散。

\noindent\textbf{链接}：\url{https://ncatlab.org/nlab/show/Nahm%27s+equation}
\end{kcard}

\begin{kcard}{注释 23：Bogomolny 方程（Bogomolny equation）}
\anchor{nLab: Bogomolny+equation}

\textbf{一句话}：Bogomolny 方程是把自对偶 Yang--Mills 方程在时间无关情形下降维后得到的一阶方程 $B_i=D_i\phi$；它的一阶解自动满足二阶运动方程，并且恰好饱和 BPS 能量界。

\textbf{数学要点}：4d 自对偶 Yang--Mills 方程是 $F=*F$。若场与时间无关、并把 $A_0$ 记作标量 $\phi$，则方程约化为 3d 的 Bogomolny 方程
\[
F_{ij}=\epsilon_{ijk}\,D_k\phi ,
\]
其中 $F_{ij}$ 是 3d 磁场、$D$ 是规范协变导数；等价地写成矢量形式 $\mathbf B=\mathbf D\phi$。单极子的能量可写成
\[
E=\int\Bigl(|\mathbf B-\mathbf D\phi|^2+2\,\mathbf B\cdot\mathbf D\phi\Bigr)
=\int|\mathbf B-\mathbf D\phi|^2+4\pi v|n| ,
\]
第二项是拓扑荷（磁荷 $n$ 与真空期望 $v$），因此 $E\ge 4\pi v|n|$，等号当且仅当 Bogomolny 方程成立——这就是单极子谱的 BPS 界与一阶方程的关系。Atiyah--Hitchin 的专著《The geometry and dynamics of magnetic monopoles》（1988）把方程列为 (2.1)，并系统研究其模空间的超凯勒度量。K--W 第 10 节正是用带边界条件的 Bogomolny 方程来刻画 Hecke 修改空间。

\textbf{在 K--W 中的位置}：§10 的标题就是 Bogomolny 方程与 Hecke 修改空间（N10、N9、M3）——'t~Hooft 算子附近的 BPS 方程加上边界条件，其解空间就是 Hecke 修改空间；§9 的 't~Hooft 算子与 §11 的膜构造都以此为基础。

\textbf{读原文前需要的前置}：注释 19、注释 22。

\textbf{常见坑}：Bogomolny 是一阶方程，只描述 BPS 扇区，不是全部运动方程；4d 瞬子与 3d 单极子是不同维数的对象，别把两者的模空间混为一谈；从自对偶方程降到 Bogomolny 方程需要“时间无关 + $A_0=\phi$”的规范选择，不是每个解都能这样写。

\noindent\textbf{链接}：\url{https://ncatlab.org/nlab/show/Bogomolny+equation}
\end{kcard}

\begin{kcard}{注释 24：$SL(2,\ZZ)$ 与模形式（the modular group and modular forms）}
\anchor{nLab: modular+group}

\textbf{一句话}：$\mathrm{SL}(2,\ZZ)$ 是 $N=4$ 超杨--米尔斯的 S-对偶群，它以 $\tau\mapsto(a\tau+b)/(c\tau+d)$ 作用在复化耦合 $\tau=\theta/2\pi+4\pi i/g^2$ 上；模形式是按这个作用、带权重因子变换的全纯函数，因此是表达 S-对偶不变量最自然的语言。

\textbf{数学要点}：$\mathrm{SL}(2,\ZZ)$ 由 $T:\tau\mapsto\tau+1$ 与 $S:\tau\mapsto-1/\tau$ 生成，其中心 $\{\pm1\}$ 的商 $\mathrm{PSL}(2,\ZZ)$ 是环面微分同胚类的映射类群，作用在复上半平面 $\mathbb H$ 上有标准基本域。权为 $k$ 的模形式是全纯函数 $f:\mathbb H\to\CC$，满足
\[
f\!\left(\frac{a\tau+b}{c\tau+d}\right)=(c\tau+d)^k f(\tau),
\qquad \begin{pmatrix}a&b\\ c&d\end{pmatrix}\in \mathrm{SL}(2,\ZZ),
\]
并在尖点处全纯（尖点形式还要求消失）；Eisenstein 级数、$j$ 不变量与 $q=e^{2\pi i\tau}$ 展开是基本例子。在 $N=4$ 理论中，Montonen--Olive 对偶把 $\tau\mapsto-1/\tau$（即 $S$）与强弱耦合互换联系起来，整群 $\mathrm{SL}(2,\ZZ)$ 是 S-对偶群；K--W 的要点 M2 说的正是这个 $S$ 变换在 2d 侧变成镜像对称。

\textbf{在 K--W 中的位置}：§2 回顾 $N=4$ 的 S-对偶与 $\tau$ 的作用（N2、S1）；§4--§5 的 S-对偶=镜像对称（M2、S4）就是 $S:\tau\mapsto-1/\tau$ 的几何化；附录对照 CFT 路线时，模形式与顶点算子代数的模不变性（配分函数在模群下的变换）是同一套语言。

\textbf{读原文前需要的前置}：注释 18、注释 14。

\textbf{常见坑}：$\mathrm{SL}$ 与 $\mathrm{PSL}$ 差一个 $\pm1$，谈“模群作用”时先确认用的是哪一个；权重约定有 $k$ 与 $2k$ 两套，变换因子是 $(c\tau+d)^k$ 还是其平方要核对；只有单连通或特定类型的规范群 S-对偶群才恰好是 $\mathrm{SL}(2,\ZZ)$，其他情形是子群或算术格，不要一律照搬。

\noindent\textbf{链接}：\url{https://ncatlab.org/nlab/show/modular+group}
\end{kcard}

\subsection*{第三组\quad 后半程代数几何：Hecke、D-模与几何 Langlands}

\begin{kcard}{注释 25：相干层（coherent sheaf）}
\anchor{nLab: coherent+sheaf}

\textbf{一句话}：相干层是环空间上“局部由有限个截生成、且任意有限个截之间的核仍有限生成”的 $\mathcal O$-模层；它是有限生成模的几何化，也是 B-模型开弦对象（B-膜）的基本材质。

\textbf{数学要点}：$\mathcal O_X$-模层 $\mathcal E$ 称为有限生成（有限型），若每点有邻域使 $\mathcal O^n|_U\twoheadrightarrow\mathcal E|_U$；称为相干，若它有限生成、且任意态射 $\mathcal O^p|_U\to\mathcal E|_U$ 的核仍有限生成；称为有限展示，若有正合列 $\mathcal O^p\to\mathcal O^n\to\mathcal E\to0$（$p,n$ 有限）。在 Noether 概形上，相干与有限展示等价；一般概形上则不等价，连 $\mathcal O_X$ 本身都可能是反例（EGA/FAC 与 Hartshorne 的定义差别即在此）。Serre 的 FAC 定理把射影簇 $\mathrm{Proj}\,R$ 上的相干层范畴等同于有限生成分次 $R$-模范畴模去有限长度子范畴；同一条短正合列中两层相干则第三层也相干。对相干层可定义秩、次数与斜率，从而谈稳定相干层，其模空间是 Donaldson--Thomas 理论与 B-模型计数的舞台；其导出范畴记作 $D^b(\mathrm{Coh}\,X)$。

\textbf{在 K--W 中的位置}：§5 的广义 B-模型即广义复几何（N5、S4），其中的 B-膜就是 $X$ 上的相干层或其复形；§12 的三种膜里 $(B,B,B)$ 型膜以相干层为对象（N12）。Hitchin 模空间 $\MH$ 的 B-膜导出范畴是整个 S-对偶/几何 Langlands 故事的“B 侧”。

\textbf{读原文前需要的前置}：注释 13、注释 15。

\textbf{常见坑}：相干层不等于向量丛（局部自由），有挠的相干层很常见，别把“相干”当“局部自由”；“有限生成”与“相干”在非 Noether 情形下不同，换定义前先看作者用哪套；不同文献的“coherent”约定（EGA 与 Hartshorne）不一致，比较定理时要核对。

\noindent\textbf{链接}：\url{https://ncatlab.org/nlab/show/coherent+sheaf}
\end{kcard}

\begin{kcard}{注释 26：导出范畴（derived category）}
\anchor{nLab: derived+category}

\textbf{一句话}：导出范畴是把链复形范畴沿拟同构局部化后得到的三角范畴，是同调代数“正确的工作场所”，使导出函子（$\mathrm{Ext}$、$\mathrm{Tor}$、$\mathrm{RHom}$、$\mathrm{R}f_*$）成为真正的函子。

\textbf{数学要点}：设 $\mathcal A$ 为 abelian 范畴，$K(\mathcal A)$ 是链复形模去链同伦的范畴，把拟同构（诱导同调同构的链映射）取为弱等价，其局部化就是导出范畴 $D(\mathcal A)$；等价地，$D(\mathcal A)\simeq K(\mathcal A)/N(\mathcal A)$，其中 $N(\mathcal A)$ 是零调复形构成的零系统（null system）。$D(\mathcal A)$ 是三角范畴，带有平移函子与 distinguished triangle；有足够内射或投射对象时，可用内射/投射分辨率计算 Hom。其增强版本（dg-范畴或稳定 $(\infty,1)$-范畴）保留了三角范畴丢掉的高阶信息。导出范畴由 Verdier 在 Grothendieck 指导下引入，最初用于把 Serre 对偶推广到相对情形；对概形 $X$ 取 $D^b(\mathrm{Coh}\,X)$，就是 B-模型与 Fourier--Mukai 变换的天然舞台。

\textbf{在 K--W 中的位置}：§5 的广义 B-模型（N5、S4）中，B-膜构成相干层的导出范畴；HMS 的表述 $\mathrm{Fuk}(X)\simeq D^b(\mathrm{Coh}\,\hat X)$ 右边正是导出范畴（M2 附近的镜像对称讨论）；§12 的 $(B,B,B)$ 膜也是以导出范畴为对象（N12）。

\textbf{读原文前需要的前置}：注释 25、注释 15。

\textbf{常见坑}：导出范畴不是 abelian 范畴，只有三角结构，不能在里面随意取核与余核；拟同构不是同构，$D(\mathcal A)$ 里才有逆；只知道三角范畴不够——增强（dg/$A_\infty$/稳定 $\infty$）信息在三角化时丢失，涉及函子性时要留意。

\noindent\textbf{链接}：\url{https://ncatlab.org/nlab/show/derived+category}
\end{kcard}

\begin{kcard}{注释 27：Fourier--Mukai 变换（Fourier--Mukai transform）}
\anchor{nLab: Fourier-Mukai+transform}

\textbf{一句话}：Fourier--Mukai 变换是以 $X\times Y$ 上的一个复形 $E$ 为核、把导出范畴 $\mathrm D(X)\to\mathrm D(Y)$ 的“拉回---张量---推前”三步复合起来的范畴化积分变换；它把“线性映射由矩阵表示”这一事实提升到导出范畴层面。

\textbf{数学要点}：设 $p,q:X\times Y\to X,Y$ 为两个投影，核 $E\in\mathrm D(X\times Y)$，则
\[
\Phi_E:\mathrm D(X)\longrightarrow \mathrm D(Y),\qquad
A\longmapsto \mathbf R q_*\!\left(\mathbf L p^*A\otimes^{\mathbf L}E\right),
\]
（所有操作都取导出函子）。原始情形（Mukai 1981）取 $X$ 为阿贝尔簇、$Y=X^\vee$ 为对偶阿贝尔簇、$E$ 为 Poincaré 线丛，得到 $\mathrm D(X)\simeq\mathrm D(X^\vee)$；一般情形下它是核积分的范畴化，对偶于“函数由积分核表示”。Orlov 定理：光滑射影簇之间任一三角全忠实函子都是 Fourier--Mukai 变换，且核在 isomorphism 意义下唯一。在增强层面有 $\mathrm{Func}^L(\mathrm D(X),\mathrm D(Y))\simeq \mathrm D(X\times Y)$。若 $X$ 是曲线上的线丛模空间，则 FM 变换的变体给出阿贝尔情形的几何 Langlands 对应（Frenkel 2005 §4.4--4.5）；T-对偶下 RR 场的变换也呈 FM 型（Hori 公式）。

\textbf{在 K--W 中的位置}：§5 的 S-对偶=镜像对称（M2、S4）在导出范畴层面的精确表述就是 Fourier--Mukai 型变换；HMS 的 $\mathrm{Fuk}(X)\simeq \mathrm D^b(\mathrm{Coh}\,\hat X)$ 两边比较时，FM 核是构造等价的标准工具；§9--§10 的 Hecke 对应与 Hecke 特征层也由 FM 型算子生成（N9、S7）。

\textbf{读原文前需要的前置}：注释 25、注释 26。

\textbf{常见坑}：Fourier--Mukai 变换一般不是等价，只有核满足特殊条件（如 Poincaré 丛、或者与 HMS 匹配的核）时才给出等价；Orlov 定理要求函子全忠实，不是“所有三角函子都来自 FM”（反例见 Rizzardo--Van den Bergh）；三角范畴层面的 FM 与稳定 $\infty$-范畴层面的 FM 在函子性上有差别，比较定理时要说明增强。

\noindent\textbf{链接}：\url{https://ncatlab.org/nlab/show/Fourier-Mukai+transform}
\end{kcard}

\begin{kcard}{注释 28：$D$-模（D-module）}
\anchor{nLab: D-module}

\textbf{一句话}：$D$-模是微分算子层 $D_X$ 上的（拟相干）模层，等价地是“带平坦联络的向量丛”的推广；它把线性偏微分方程系统放进同调代数与层论的框架。

\textbf{数学要点}：$D_X$ 是由向量场生成的微分算子层，$D_X$-模自动是 $\mathcal O_X$-模（因为 $\mathcal O_X\subset D_X$ 是零阶算子），但反向不成立。在光滑复流形上，$D_X$-模范畴等价于 de Rham 空间 $\mathrm{dR}(X)$（de Rham stack）上的拟相干层范畴，这解释了“$D$-模=带平坦联络的模”这一直观。若一个 $D$-模作为 $\mathcal O_X$-模相干，它同构于某个全纯向量丛带平坦联络的截面层；进一步，作为 $D_X$-模相干的（nonsingular）$D$-模构成的子范畴等价于局部系统范畴（Koszul--Malgrange）。$D$-模有完整的六函子形式（拉回、推前、张量、内 Hom 等），Riemann--Hilbert 对应则说正则完整的全纯 $D$-模对应于可构造层。$D$-模是几何表示论（Beilinson--Bernstein 局部化）与几何 Langlands 自守侧的核心语言：Hecke 特征层正是 $D$-模。

\textbf{在 K--W 中的位置}：§11 的 A-膜与 $D$-模对应（N11、S8、M5）：$\MH$ 上的典范余各向同性 A-膜 $\mathcal A_{cc}$ 通过 $\mathrm{Hom}(\mathcal A_{cc},-)$ 的 sheafification 给出 $D$-模（注释 16）；§12 用三种膜把规范理论对象翻译成 $D$-模（N12）；§1 引言的第三条数学史线索就是 $D$-模与几何 Langlands。

\textbf{读原文前需要的前置}：注释 11、注释 12。

\textbf{常见坑}：$D$-模不是“向量丛”本身，$\mathcal O_X$-相干不等于 $D_X$-相干（正特征情形尤其明显）；Riemann--Hilbert 对应只对正则完整（regular holonomic）$D$-模成立，别对任意 $D$-模谈 RH；左、右 $D$-模的约定不同，写公式前先核对作者用哪种。

\noindent\textbf{链接}：\url{https://ncatlab.org/nlab/show/D-module}
\end{kcard}

\begin{kcard}{注释 29：仿射 Grassmannian（affine Grassmannian）}
\anchor{nLab: affine+Grassmannian}

\textbf{一句话}：仿射 Grassmannian $\mathrm{Gr}_G$ 是 loop group $G(k((t)))$ 的“旗流形”，它参数化形式圆盘上的 $G$-丛及其在穿孔圆盘上的平凡化，是几何 Satake 等价的几何舞台。

\textbf{数学要点}：$\mathrm{GL}_n$ 情形：对 $k$-代数 $R$，$\mathrm{Gr}_{\mathrm{GL}_n}(R)$ 是 $R((t))^n$ 中的格（lattice）$\Lambda$ 的集合，即有限生成投射子模且 $\Lambda\otimes_{k[[t]]}k((t))=k((t))^n$。一般约化群 $G$：令 $D_R=\operatorname{Spec} k[[t]]\ \widehat\times\operatorname{Spec} R$、$D^*_R=\operatorname{Spec} k((t))\ \widehat\times\operatorname{Spec} R$，$\mathrm{Gr}_G(R)$ 是数对 $(\mathcal E,\beta)$ 的集合，其中 $\mathcal E$ 是 $D_R$ 上的 $G$-丛、$\beta:\mathcal E|_{D^*_R}\xrightarrow{\ \simeq\ }\mathcal E^0|_{D^*_R}$ 是到平凡丛的平凡化。它是 ind-scheme（可看作 $G((t))/G[[t]]$ 的 ind-商），作为函子 ind-representable；$\mathrm{SL}_n$ 情形可嵌入 Sato Grassmannian。它与 Ran 空间、几何 Langlands 直接相关；Zhu 的讲义《An introduction to affine Grassmannians and the geometric Satake equivalence》（arXiv:1603.05593）是标准入口。

\textbf{在 K--W 中的位置}：§9 的 't Hooft 算子与 Hecke 修改（N9、S7）：'t Hooft 算子插入点附近的 Hecke 修改是纯粹局部的，其参数空间正是 $\mathrm{Gr}_G$（Schubert 胞腔分解给出荷的格）；§10 的 Bogomolny 方程与“Hecke 修改空间”讨论单极子冒泡与边界条件时同样落在 $\mathrm{Gr}_G$ 上；注释 32 的几何 Satake 等价就建立在 $\mathrm{Gr}_G$ 的等变 perverse 层范畴上。

\textbf{读原文前需要的前置}：注释 1、注释 30。

\textbf{常见坑}：$\mathrm{Gr}_G$ 是 ind-scheme 而非有限维簇，谈“维数”只能逐个 Schubert 胞腔谈（每个有限维）；把 $\mathrm{Gr}_G$ 写成 $G((t))/G[[t]]$ 时是 ind-商，拓扑（étale/fppf）选择要说明；不同作者用“格”与“平凡化对”两种等价定义，读公式前先对齐。

\noindent\textbf{链接}：\url{https://ncatlab.org/nlab/show/affine+Grassmannian}
\end{kcard}

\begin{kcard}{注释 30：Hecke 对应与 Hecke 修正（Hecke correspondence and modification）}
\anchor{nLab: Hecke+correspondence}

\textbf{一句话}：Hecke 对应是丛模空间 $\mathrm{Bun}_G$ 到自身的 correspondence（Hecke 栈 $\mathrm{Hck}$ 上的两个投影），在一点处对丛做“修正”（改变该点附近局部结构）给出作用在层上的 Hecke 变换；Hecke 特征层就是它的本征对象。

\textbf{数学要点}：Hecke 栈 $\mathrm{Hck}$ 参数化三元组 $(\mathcal E_1,\mathcal E_2,\beta)$，其中 $\beta:\mathcal E_1\simeq\mathcal E_2$ 在除有限个点外是同构（这正是“修正”/modification）：在曲线上一点处，$\beta$ 允许把 $\mathcal E_1$ 的局部结构换成另一格子。于是有 correspondence
\[
\mathrm{Bun}_G\xleftarrow{\ \pi_1\ }\mathrm{Hck}\xrightarrow{\ \pi_2\ }\mathrm{Bun}_G,
\]
Hecke 变换形如 $\Phi(A)=\pi_{2*}\big(\pi_1^*A\otimes\rho\big)$。对曲线 $X$ 上的 $\mathrm{Bun}_G$，这与经典模曲线情形 $X_0(p)\rightrightarrows X$ 及 Hecke 算子 $T_p$ 作用在模形式 $H^0(X(\Gamma),\omega^k)$ 上完全平行：经典 Hecke 算子就是“线化后的 Hecke 对应”。对 $\mathrm{Bun}_G$ 的 Hecke 对应是几何 Langlands 的“算子”侧；nLab 页给出与 S-对偶的对照表：Hecke 变换 $\leftrightarrow$ 't Hooft 算子，Hecke 特征层 $\leftrightarrow$ 磁特征膜（magnetic eigenbrane）。标准参考是 Frenkel 2005 §3.7。

\textbf{在 K--W 中的位置}：§9 't Hooft 算子与 Hecke 修改（N9、S7、M3：“'t Hooft 算子作用 = Hecke 修改”）；§10 Bogomolny 方程与 Hecke 修改空间（边界条件、单极子冒泡给出 $\mathrm{Gr}_G$ 上的修正）；§3 扭化后 't Hooft 线算子在二维视图中也通过 Hecke 对应编码。

\textbf{读原文前需要的前置}：注释 1、注释 29。

\textbf{常见坑}：Hecke 对应不是自同构，而是多值 correspondence，别把它当成 $\mathrm{Bun}_G$ 的端同构；修正只发生在有限个点，其余处是同构；“Hecke 栈”与“Hecke 范畴”（几何 Satake 的等变层范畴）是不同对象，注意区分。

\noindent\textbf{链接}：\url{https://ncatlab.org/nlab/show/Hecke+correspondence}
\end{kcard}

\begin{kcard}{注释 31：Hecke 特征层（Hecke eigensheaf）}
\anchor{nLab: Hecke+eigensheaf}

\textbf{一句话}：Hecke 特征层是在 Hecke 变换下变成自身的张量倍数（$\Phi(A)\cong A\otimes V$）的对象——它是几何 Langlands 预言的“本征函数”在层论中的化身：给定曲线上的局部系统 $E$，应当存在 $\mathrm{Bun}_G$ 上的层 $A_E$，它在所有 Hecke 变换下本征。

\textbf{数学要点}：Hecke 对应给出 $\mathrm{Bun}_G\xleftarrow{\pi_1}\mathrm{Hck}\xrightarrow{\pi_2}\mathrm{Bun}_G$ 与 Hecke 变换（注释 30）。一个对象 $A$ 若满足
\[
\Phi^{\,V}(A)\ \cong\ A\otimes V
\]
对每个修正/每个 $V$ 成立，就称为 Hecke 特征层。几何 Langlands 猜想（对光滑射影曲线 $X$ 上的 $G$-丛）断言：对 $X$ 上每个不可约 $G$-局部系统 $E$（等价地 $\Ldual$-局部系统），存在 $\mathrm{Bun}_G$ 上的 $D$-模（或 perverse sheaf）$\mathcal A_E$，使得对每个 Hecke 修正它本征，且本征值由几何 Satake 等价（注释 32）识别为 $\mathrm{Rep}(\Ldual)$ 中的表示。因此 Hecke 特征层把“Langlands 对偶群的表示”组织成 $\mathrm{Bun}_G$ 上层的本征值。它与 S-对偶中的磁特征膜（'t Hooft 算子的本征膜）对应；标准参考 Frenkel 2005 §3.8。其构造与存在性是几何 Langlands 的核心问题。

\textbf{在 K--W 中的位置}：§9--§10 的 't Hooft 算子与 Hecke 修改（N9、S7、M3）；§11 的 A-膜 $\to$ $D$-模对应（N11、S8、M5）最终目标的输出就是 Hecke 特征层；§4--§5 的 Kapustin--Witten TQFT 预言整个流程（电特征膜 $\to$ S-对偶 $\to$ 磁特征膜 $\to$ $D$-模）的终点。

\textbf{读原文前需要的前置}：注释 28、注释 30。

\textbf{常见坑}：特征层不是“特征向量”，而是随 Hecke 算子变化的一整族层，本征值本身是表示论数据；“Hecke 特征层”有 perverse sheaf 与 $D$-模两种语言，比较时要用 RH 对应；本征值必须是“up to tensor multiple”，别写成严格同构。

\noindent\textbf{链接}：\url{https://ncatlab.org/nlab/show/Hecke+eigensheaf}
\end{kcard}

\begin{kcard}{注释 32：几何 Satake 等价（geometric Satake equivalence）}
\anchor{nLab: geometric+Satake+equivalence}

\textbf{一句话}：几何 Satake 等价把仿射 Grassmannian 上等变 $D$-模的卷积范畴等同于 Langlands 对偶群的表示范畴，是经典 Satake 同构的范畴化。

\textbf{数学要点}：经典 Satake 同构：球 Hecke 代数 $\mathcal H$ 是双陪集空间 $G(\mathcal O)\backslash G(K)/G(\mathcal O)$ 上紧支撑 $\ZZ$-值函数在卷积下的环，Gelfand 技巧说明它交换；则
\[
\mathcal H\otimes\ZZ[q^{\pm 1/2}]\ \simeq\ R(\Ldual)\otimes\ZZ[q^{\pm 1/2}],
\]
其中 $q$ 是 $\mathcal O$ 剩余域的大小，$R(\Ldual)=K_0(\mathrm{Rep}(\Ldual))$ 是对偶群的表示环。几何版本：$G$ 为 $\CC$ 上约化群、$\Ldual$ 为其对偶群，$\mathcal L G$ 为 loop group、$\mathcal L^+G$ 为 arc group、$\mathrm{Gr}_G=\mathcal L G/\mathcal L^+G$；则
\[
\mathrm{DMod}(\mathrm{Gr}_G)^{\mathcal L^+G}\ \simeq\ \mathrm{Rep}(\Ldual)
\]
是配卷积 monoidal 结构的对称 monoidal 阿贝尔范畴等价。这里交换性约束（commutativity constraint）是额外汇编的数据，而不是 monoidal 结构自动给出的性质。几何 Satake 是几何 Langlands 的“局部”一半：Hecke 特征层的本征值属于 $\mathrm{Rep}(\Ldual)$。

\textbf{在 K--W 中的位置}：§9 的 't Hooft 算子与 Hecke 修正（N9）背后的表示论机器；§10 的单极子冒泡与 Hecke 修改空间通过 $\mathrm{Gr}_G$ 的 Schubert 胞腔与 Satake 参数对应；§11 的 A-膜 $\to$ $D$-模构造需要局部 Satake 来识别本征值；附录对照 VOA/CFT 路线时（注释 36），几何 Satake 是代数侧的支柱。

\textbf{读原文前需要的前置}：注释 1、注释 29、注释 34。

\textbf{常见坑}：等变群是 arc group $\mathcal L^+G$，不是 $G$ 本身；交换性约束是额外数据，不能默认唯一；$D$-模版本与 perverse sheaf 版本通过 RH 对应比较，别混用。

\noindent\textbf{链接}：\url{https://ncatlab.org/nlab/show/geometric+Satake+equivalence}
\end{kcard}

\begin{kcard}{注释 33：几何 Langlands 纲领（geometric Langlands correspondence）}
\anchor{nLab: geometric+Langlands+correspondence}

\textbf{一句话}：几何 Langlands 对应是数论 Langlands 对应在“函数域/复曲线”上的类比：它猜想 $\mathrm{Bun}_G(\Sigma)$ 上的 $D$-模导出范畴与 $\mathrm{Loc}_{\Ldual}(\Sigma)$ 上的拟相干层范畴等价，且这个等价是非阿贝尔版本的 Fourier--Mukai 积分变换。

\textbf{数学要点}：通过函数域类比（数域 $\to$ 函数域 $\to$ 复曲线上的有理函数，Galois 群 $\to$ 基本群，Galois 表示 $\to$ 平坦联络/局部系统，自守表示 $\to$ $\mathrm{Bun}_G$ 上的 $D$-模），Beilinson--Drinfeld 与 Frenkel 2005 的原始猜想是导出范畴等价
\[
\mathcal O\mathrm{Mod}\big(\mathrm{Loc}_{\Ldual}(\Sigma)\big)\ \simeq\ \mathcal D\mathrm{Mod}\big(\mathrm{Bun}_G(\Sigma)\big),
\]
且该等价把左边的 skyscraper sheaves（范畴化的 Dirac 分布）送到右边的 Hecke 特征层（注释 31）。$G$ 为环面（阿贝尔情形）时等价已证，且恰由 Fourier--Mukai 变换给出（Laumon 85/96、Rothstein 96）。朴素形式对一般 $G$ 是错的：$G=\mathrm{SL}_2$、$\Sigma=\mathbb P^1$ 时失败（Lafforgue 09）；Arinkin--Gaitsgory 2012 用拟相干层的“幂零”ind-对象与稳定 $\infty$-范畴修正为 categorical geometric Langlands 猜想，nLab 记载该形式已由 Gaitsgory 等（2025）证明。物理上，$\mathrm{Bun}_G$ 上的 $D$-模（Hitchin 联络、共形块的丛）正是 Chern--Simons 理论全息对偶 WZW 模型时出现的东西；Kapustin--Witten 2006 提出几何 Langlands 是 S-对偶/镜像对称的特例，但这一物理推导尚未变成定理。

\textbf{在 K--W 中的位置}：§2 的 S-对偶与对偶群同构（S1、N2）；§4--§5 的 $\MH(G)$ 与 $\MH(\Ldual)$ 的镜像对称（S3、S4、M2）；§9--§10 的 Hecke 对应（S7、N9）；§11 的 $D$-模构造（S8、N11、M5）；§12 用三种膜把这些对应物化（N12）。

\textbf{读原文前需要的前置}：注释 28、注释 31、注释 34。

\textbf{常见坑}：朴素形式（普通拟相干层 vs $D$-模）对一般 $G$ 是错的，必须用 Arinkin--Gaitsgory 的修正形式；“几何 Langlands”有拟相干层、ind-幂零、扭化 $D$-模等多个版本，别混用；“物理实现”是启发式推导，不等同于定理证明。

\noindent\textbf{链接}：\url{https://ncatlab.org/nlab/show/geometric+Langlands+correspondence}
\end{kcard}

\begin{kcard}{注释 34：Langlands 对偶群（Langlands dual group）}
\anchor{nLab: Langlands+program}

\textbf{一句话}：Langlands 对偶群 $\Ldual$ 是把约化群 $G$ 的根与余根（对偶根数据）互换后得到的群；它把 $G$ 的表示论与 $G$ 的 $L$-函数联系起来，在几何 Langlands 中它就是“镜像”的另一侧。

\textbf{数学要点}：对偶根数据把 $G$ 的根格与余根格互换：$X^*(\Ldual)=X_*(G)$、$X_*(\Ldual)=X^*(G)$；非分裂情形要取 $L$-群 $\Ldual=G^\vee\rtimes\mathrm{Gal}$，而不只是对偶群。典型例子：$\mathrm{GL}_n$ 自对偶，$\mathrm{SL}_n\leftrightarrow\mathrm{PGL}_n$，$\mathrm{Sp}(2n)\leftrightarrow\mathrm{SO}(2n+1)$，$G_2$ 自对偶。它在数论 Langlands 中的作用是组织 $L$-函数：对约化群 $G$，$L$-函数要附加一个对偶群表示 $\Ldual\to\mathrm{GL}_n$ 的数据；nLab 的 Langlands program 页列出三条猜想——$L$-函数的解析延拓与函数方程、以及函子性（$\Ldual\to\Ldual{}'$ 的 homomorphism 对应自守表示的移动）；在几何情形，Galois 表示的角色由曲线上的平坦联络/局部系统扮演（注释 11），自守表示的角色由 $\mathrm{Bun}_G$ 上的 $D$-模扮演（注释 28）。

\textbf{在 K--W 中的位置}：§2 的 $S$-对偶 $\tau\to-1/\tau$ 与对偶群同构 $G\leftrightarrow\Ldual$（S1、N2，注释 18）；§3 的 't Hooft 荷与 $\tau$ 的作用（N3）；§4 的 $\MH(G)$ 与 $\MH(\Ldual)$ 互为镜像（S3、N3）；§9 的 Hecke 特征层本征值落在 $\mathrm{Rep}(\Ldual)$（S7，注释 31、32）；§12 的三种膜在 $G$ 与 $\Ldual$ 两侧互换（N12）。

\textbf{读原文前需要的前置}：注释 1、注释 35。

\textbf{常见坑}：对偶群一般不同构于 $G$，只有自对偶情形（$\mathrm{GL}_n$、$G_2$ 等）才同构；非分裂群要用 $L$-群而非对偶群；Langlands 本人没用“dual group”一词，原始文献称 $L$-group（nLab 页有这条术语注记）。

\noindent\textbf{链接}：\url{https://ncatlab.org/nlab/show/Langlands+program}
\end{kcard}

\begin{kcard}{注释 35：根格与权格（root lattice and weight lattice）}
\anchor{nLab: root in representation theory}

\textbf{一句话}：根格/余根格与权格/余权格是约化群与李代数的组合骨架：根系决定群（连同中心），权格是根格的对偶格，表示的最高权落在权格中，而 Langlands 对偶正是把这两套格互换。

\textbf{数学要点}：根系（nLab “root”）：在欧氏空间 $E$ 中，有限非零向量集 $\Delta$ 若张成 $E$、每个反射 $\sigma_\alpha$ 置换 $\Delta$、且整性条件 $\alpha\dashv\beta:=2\langle\alpha|\beta\rangle/\langle\beta|\beta\rangle\in\ZZ$ 成立，则称 $\Delta$ 为根系；若 $\operatorname{span}(\alpha)\cap\Delta=\{-\alpha,\alpha\}$ 则称 reduced。Weyl 群 $W(\Delta)$ 由反射 $\sigma_\alpha$ 生成。权（nLab “weight”）：紧连通 Lie 群 $G$ 的极大环面 $T$ 的不可约表示称为权；表示 $V$ 中按权 $\alpha$ 分解的子空间是权空间（$T$-作用的 eigenspace）。对偶格：根格 $Q=\bigoplus_\alpha\ZZ\alpha$、权格 $P=\{\lambda:\langle\lambda,\alpha^\vee\rangle\in\ZZ\ \forall\alpha\}$，有 $Q\subset P$，商 $P/Q$ 对偶于中心；对偶根数据交换 $Q\leftrightarrow P^\vee$。最高权表示由 $\lambda\in P_+$ 分类，Weyl 特征公式给出特征标。这些格在规范理论里就是电荷/磁荷的量子化格。

\textbf{在 K--W 中的位置}：§2 的对偶群同构（S1）本质就是根数据对偶；§7 的 GNO 线/'t Hooft 荷与 flux 修正（N7）中，Wilson 线的荷在权格、't Hooft 线的荷在余权格（磁荷格），这正是 S-对偶 $G\leftrightarrow\Ldual$ 交换两格的内容；§9 的 't Hooft 算子与 Hecke 修正（N9）、§10 单极子与 Bogomolny 方程（N10、注释 22、23）里的荷同样由这些格量子化。

\textbf{读原文前需要的前置}：注释 1、注释 18。

\textbf{常见坑}：根格 $Q$ 与权格 $P$ 不同（$P/Q$ 反映中心），写公式别混；根 $\alpha$ 与余根 $\alpha^\vee$ 要分清；$\mathrm{SL}_n$ 与 $\mathrm{PGL}_n$ 的权格/根格正好互换，这是理解对偶群的标准例子。

\noindent\textbf{链接}：\url{https://ncatlab.org/nlab/show/root+%28in+representation+theory%29}；\url{https://ncatlab.org/nlab/show/weight+%28in+representation+theory%29}
\end{kcard}

\subsection*{第四组\quad 附录对照：CFT 路线与概念地图}

\begin{kcard}{注释 36：顶点算子代数与 CFT 路线（vertex operator algebra）}
\anchor{nLab: vertex+operator+algebra}

\textbf{一句话}：顶点算子代数（VOA）是“乘法依赖复平面上点”的代数，它编码二维共形场论的局部（genus 0）手征一半；在几何 Langlands 的 CFT 路线里，它取代 $D$-模成为算子侧/自守侧的候选语言。

\textbf{数学要点}：VOA 的数据是：非负分次复向量空间 $V=\bigoplus_{n\ge0}V_n$、真空向量 $|0\rangle\in V_0$、平移算子 $T:V\to V$、以及态--场对应 $Y(-,z):V\to\operatorname{End}V[[z,z^{-1}]]$，满足真空、平移、局域性三条公理；当 $T$ 来自“能量--动量张量”$\omega$、其 Fourier 模 $L_i$ 给出 Virasoro 代数作用时称共形顶点代数。一个 VOA 只编码手征/全纯的一半，完整 genus-0 共形场论由两个 VOA（full field algebra）拼成。Huang 的算子化刻画：VOA 恰是全纯 punctured sphere operad 的代数，即幺半函子 $2\mathrm{Cob}^0_{\mathrm{conf}}\to\mathrm{Vect}$，这把它与二维 CFT 的 $n$-点函数显式联系起来。有理 VOA（每个模完全可约）的模范畴是模张量范畴（Huang），这是 CFT 路线的核心结构；另有与 $A_\infty$-代数（Zeitlin）、factorization algebra（Gwilliam）的对应。Beilinson--Drinfeld 的 chiral algebra 是其代数几何版本，标准参考 Frenkel--Ben-Zvi《Vertex algebras and algebraic curves》。

\textbf{在 K--W 中的位置}：本卡属“附录对照”组，不参与主链 S1--S8，只提示原文为何选择层论/$D$-模路线：注释 33 所引 nLab 页把“共形块丛上的 Hitchin 联络 $D$-模”列为自守表示的几何对应物——共形块正是 VOA/仿射 Lie 代数 $\hat{\mathfrak g}$ 的模空间上的丛；而 WZW 模型的 VOA 与 §2 的规范群 $G$、§9 的仿射 Grassmannian（loop group 的商）共用同一套组合数据。因此 CFT 路线与 $D$-模路线在物理上同源，在数学上一个走表示论、一个走层论。

\textbf{读原文前需要的前置}：注释 28、注释 29、注释 32。

\textbf{常见坑}：VOA 不是普通结合代数（乘法依赖点，且以 OPE 形式出现）；它只编码手征一半，别把它等同于整个二维 CFT；“有理（rational）”有严格定义（所有模完全可约），不是日常意义的“有理性”。

\noindent\textbf{链接}：\url{https://ncatlab.org/nlab/show/vertex+operator+algebra}
\end{kcard}

\subsection*{收尾：概念地图与三条阅读路线}

把 36 张卡按依赖关系摆在一起，可以得到下面这张地图（箭头表示“读后者之前最好先读前者”，不是数学蕴含）：

\begin{center}
\begin{tikzpicture}[font=\footnotesize,
  box/.style={draw, rounded corners=2pt, align=center, inner sep=3pt,
              text width=2.6cm, minimum height=0.95cm}]
\node[box] (bun)  at (0,0)      {主丛、稳定丛\\模空间与 GIT};
\node[box] (hig)  at (3.4,0)    {Higgs 丛与\\Hitchin 方程};
\node[box] (mh)   at (6.8,0)    {$\MH$：超凯勒\\Hitchin 纤维化};
\node[box] (sym)  at (6.8,-1.7) {$N=4$ SYM 的\\拓扑扭化};
\node[box] (tft)  at (6.8,-3.4) {2d $A/B$ 模型\\镜像对称};
\node[box] (brn)  at (3.4,-3.4) {D-膜：余各向同性\\$A$-膜};
\node[box] (loc)  at (0,-3.4)   {平坦联络/\\局部系统};
\node[box] (hck)  at (3.4,-5.1) {Hecke 对应\\与 Hecke 特征层};
\node[box] (dm)   at (6.8,-5.1) {$D$-模与\\几何 Langlands};
\draw[->] (bun) -- (hig);
\draw[->] (hig) -- (mh);
\draw[->] (mh) -- (sym);
\draw[->] (sym) -- (tft);
\draw[->] (tft) -- (brn);
\draw[->] (brn) -- (hck);
\draw[->] (hck) -- (dm);
\draw[->] (loc) -- (brn);
\end{tikzpicture}
\end{center}

地图的三层恰好对应原文的叙事顺序：底层（左）是数学前置，中层（右）是物理语言，顶层（右下）是 Langlands 输出。据此给出三条阅读路线：

\begin{itemize}
\item \textbf{跟原文主线（最短闭环）}：3（超凯勒）$\to$ 4（超凯勒商）$\to$ 5（Higgs 丛）$\to$ 6（$\MH$）$\to$ 7（Hitchin 方程）$\to$ 8（Hitchin 纤维化）$\to$ 9（非阿贝尔 Hodge）$\to$ 10（广义复几何）$\to$ 13（$A/B$ 模型）$\to$ 14（镜像）$\to$ 15/16（膜）$\to$ 17（扭化）$\to$ 18（$N=4$）$\to$ 28（$D$-模）$\to$ 30/31（Hecke）$\to$ 33（几何 Langlands）。
\item \textbf{物理直觉路线}：18（$N=4$）$\to$ 17（扭化）$\to$ 13（$A/B$ 模型）$\to$ 14（镜像）$\to$ 15（膜）$\to$ 19（BPS）$\to$ 20（BRST）$\to$ 21/22/23（瞬子、单极子、Bogomolny）。
\item \textbf{代数几何路线}：25（相干层）$\to$ 26（导出范畴）$\to$ 27（Fourier--Mukai）$\to$ 28（$D$-模）$\to$ 29（仿射 Grassmannian）$\to$ 30--32（Hecke/几何 Satake）$\to$ 33（几何 Langlands）$\to$ 34/35（对偶群、根格权格）。
\end{itemize}

最后提醒两点。其一，原文的五个“miracle”（$\mathbb{CP}^1$ 族保持拓扑性、4d $S$-对偶在 2d 侧变成镜像、't Hooft 作用等于 Hecke 修改、Wilson 特征膜等于局部系统、典范余各向同性膜给出 $A$-膜到 $D$-模的桥）分别对应注释 17、14、30、15/16、16/28——若读原文时卡在某个“miracle”，先回到对应卡片。其二，26 以后各卡的链接是几何 Langlands 文献的入口，但原文 §11--§12 的论证只需要注释 28、30、31、33 四张卡的核心定义；其余卡片是给想继续深挖的读者的路标，不必在初读时全部消化。

% __NEXT__
